% ----------------------------------------------------------
\begin{frame}{Join $\rajoin$ --- Why We Need It}
  Relational databases deliberately \textbf{split} information across multiple
  tables to reduce redundancy.  A \textbf{join} reassembles related tuples from
  two (or more) tables on demand.

  \begin{block}{Core Identity}
    \[
      R_1 \;\rajoin_C\; R_2 \;=\; \rasel_C\!\bigl(R_1 \times R_2\bigr)
    \]
    A join is \emph{syntactic sugar} for a Cartesian product followed by a
    restriction.  The condition $C$ discards meaningless combinations.
  \end{block}

  \begin{examples}{Motivating Scenario}
    \begin{itemize}
      \item \texttt{Products(ProdID, Name, BrandID)} ---  holds product data
      \item \texttt{Brands(BrandID, BrandName)} --- holds brand names
      \item To display every product with its brand name, we must
            \textbf{join} the two tables on \texttt{Products.BrandID = Brands.BrandID}.
    \end{itemize}
  \end{examples}
\end{frame}

% ----------------------------------------------------------
\begin{frame}{Types of Joins --- Overview}
  There are several join variants; each handles \emph{non-matching} tuples
  differently.

  \vspace{0.4em}
  \begin{center}
    \footnotesize
    \begin{tabular}{llp{5.5cm}}
      \toprule
      \textbf{Symbol} & \textbf{Name} & \textbf{What happens to non-matching tuples} \\
      \midrule
      $\rajoin$          & Natural Join       & Discarded; join on all shared attribute names \\
      $\rajoin_C$        & Theta / Inner Join & Discarded; explicit condition $C$ \\
      $\ralsem$          & Left Semi-Join     & Right non-matches discarded; left attrs only \\
      $\rarsem$          & Right Semi-Join    & Left non-matches discarded; right attrs only \\
      $\rajoin_L$        & Left Outer Join    & Right side filled with \texttt{NULL} \\
      $\rajoin_R$        & Right Outer Join   & Left side filled with \texttt{NULL} \\
      $\rajoin_{LR}$     & Full Outer Join    & Both sides filled with \texttt{NULL} \\
      $\triangleright$   & Anti Join          & Only left tuples \emph{without} a right match \\
      \bottomrule
    \end{tabular}
  \end{center}

  \vspace{0.4em}
  The following frames cover each variant in detail.
\end{frame}

% ----------------------------------------------------------
\begin{frame}{Natural Join $\rajoin$}
  The natural join \textbf{automatically} identifies every attribute that
  appears in both relations and joins on their equality.  Shared columns appear
  only once in the result.

  \begin{block}{Properties}
    \begin{itemize}
      \item No explicit condition needed --- shared attribute names define it.
      \item Result schema = union of both schemas (duplicates removed).
      \item Can be chained: $A \rajoin X \rajoin M$ joins three tables.
    \end{itemize}
  \end{block}

  \begin{alertblock}{Danger: Silent Wrong Results}
    If two tables coincidentally share an attribute name that should
    \textbf{not} be part of the join, the natural join silently produces
    incorrect output.  Always verify which attributes are shared before using it.
  \end{alertblock}

  \begin{examples}{Example}
    \texttt{Albums(AlbumID, ArtistID, Title)} and
    \texttt{Artists(ArtistID, Name)}\\
    $\texttt{Albums} \rajoin \texttt{Artists}$ joins on \texttt{ArtistID}
    automatically.
  \end{examples}
\end{frame}

% ----------------------------------------------------------
\begin{frame}{Inner Join / Theta Join $\rajoin_C$}
  The \textbf{inner join} requires an \emph{explicitly stated} join condition.
  The general form is the \textbf{theta join} ($\theta \in \{=, <, >, \leq,
  \geq, \neq\}$); when $\theta$ is $=$ it is called an \textbf{equi-join}.

  \begin{block}{Key Properties}
    \begin{itemize}
      \item Only tuples with a matching counterpart in the other relation
            appear in the result --- non-matching tuples are discarded.
      \item Both relations' full attribute sets appear in the result
            (unlike natural join, shared column names are \emph{not} merged).
      \item Syntax: $R_1 \;\rajoin_{R_1.A \,=\, R_2.B}\; R_2$
    \end{itemize}
  \end{block}

  \begin{examples}{Example}
    \[
      \texttt{Products} \;\rajoin_{\texttt{Products.BrandID}\,=\,\texttt{Brands.BrandID}}\; \texttt{Brands}
    \]
    Products without a matching brand \emph{and} brands without any product
    are both excluded.
  \end{examples}
\end{frame}

% ----------------------------------------------------------
\begin{frame}{Inner Join --- Venn Diagram and Example Table}
  \begin{columns}[c]
    \begin{column}{0.3\textwidth}
      \centering
      \vennjoin{0}{1}{0}\\[4pt]
      \small Only the \textbf{intersection} survives.
    \end{column}
    \begin{column}{0.65\textwidth}
      \small
      \textbf{Products}
      \begin{tabular}{lll}
        \toprule
        \textbf{ProdID} & \textbf{Name} & \textbf{BrandID}\\
        \midrule
        P1 & Headphones & B1 \\
        P2 & Speaker    & B2 \\
        P3 & Cable      & \texttt{NULL} \\
        \bottomrule
      \end{tabular}

      \vspace{6pt}
      \textbf{Brands}
      \begin{tabular}{ll}
        \toprule
        \textbf{BrandID} & \textbf{BrandName}\\
        \midrule
        B1 & Sony \\
        B2 & Bose \\
        B3 & JBL  \\
        \bottomrule
      \end{tabular}

      \vspace{6pt}
      \textbf{Result} (inner join on BrandID)
      \begin{tabular}{llll}
        \toprule
        ProdID & Name & BrandID & BrandName\\
        \midrule
        P1 & Headphones & B1 & Sony \\
        P2 & Speaker    & B2 & Bose \\
        \bottomrule
      \end{tabular}
      \small P3 (no brand) and B3 (no product) are excluded.
    \end{column}
  \end{columns}
\end{frame}

% ----------------------------------------------------------
\begin{frame}{Composing Operators --- Rename + Join + Project}
  The true power of relational algebra is \textbf{composing} operators freely:
  the output of any operator is itself a relation.

  \begin{block}{Reading Strategy}
    Read relational algebra expressions \textbf{inside-out}: start from the
    innermost sub-expression and apply each outer operator in turn.
  \end{block}

  \begin{examples}{Example: Show Product Name and Brand Name Only}
    \small
    \[
      \raproj_{\texttt{ProdID},\texttt{ProductName},\texttt{BrandName}}\!\Bigl(
        \raren_{\texttt{ProductName/Name}}(\texttt{Products})
        \;\rajoin_{\texttt{BrandID}}\;
        \raren_{\texttt{BrandName/Name}}(\texttt{Brands})
      \Bigr)
    \]
    Steps: (1) rename \texttt{Name} in Products $\to$ \texttt{ProductName};
    (2) rename \texttt{Name} in Brands $\to$ \texttt{BrandName} (avoids
    ambiguity); (3) join; (4) project.
  \end{examples}
\end{frame}

% ----------------------------------------------------------
\begin{frame}{Left Semi-Join $\ralsem$ and Right Semi-Join $\rarsem$}
  A semi-join tests for the \emph{existence} of a match but does \textbf{not}
  append columns from the other side.

  \begin{block}{Left Semi-Join $\ralsem$}
    Returns all tuples of $R_1$ for which at least one matching tuple exists in
    $R_2$.  Only columns of $R_1$ are retained.\\[2pt]
    $X \;\ralsem_{X.A\,=\,Y.A}\; Y$
  \end{block}

  \begin{block}{Right Semi-Join $\rarsem$}
    Symmetric: returns tuples of $R_2$ that have a match in $R_1$, using only
    $R_2$'s columns.\\[2pt]
    $X \;\rarsem_{X.A\,=\,Y.A}\; Y$
  \end{block}

  \begin{examples}{Use Case --- ``Which customers have placed at least one order?''}
    \[
      \texttt{Customers} \;\ralsem_{\texttt{Customers.ID}\,=\,\texttt{Orders.CustomerID}}\; \texttt{Orders}
    \]
    Each customer appears \textbf{at most once} in the result (no row-per-order
    duplication); no order data is added.
  \end{examples}
\end{frame}

% ----------------------------------------------------------
\begin{frame}{Left Outer Join $\rajoin_L$ --- Definition and Venn Diagram}
  \begin{columns}[c]
    \begin{column}{0.3\textwidth}
      \centering
      \vennjoin{1}{1}{0}\\[4pt]
      \small \textbf{All} of A; matching part of B.
    \end{column}
    \begin{column}{0.65\textwidth}
      \begin{block}{Left Outer Join $\rajoin_L$}
        All tuples from the \textbf{left} relation are kept.  Matching tuples
        from the right are appended.  Where no match exists, right-side columns
        are filled with \texttt{NULL}.
      \end{block}

      \vspace{4pt}
      \small
      \textbf{Products} $\rajoin_L$ \textbf{Brands}:

      \begin{tabular}{llll}
        \toprule
        ProdID & Name & BrandID & BrandName\\
        \midrule
        P1 & Headphones & B1 & Sony \\
        P2 & Speaker    & B2 & Bose \\
        P3 & Cable      & \texttt{NULL} & \texttt{NULL} \\
        \bottomrule
      \end{tabular}

      \vspace{4pt}
      P3 is preserved; its brand columns are \texttt{NULL}.\\
      B3 (JBL, no products) is \textbf{not} included.
    \end{column}
  \end{columns}
\end{frame}

% ----------------------------------------------------------
\begin{frame}{Right Outer Join $\rajoin_R$ and Full Outer Join $\rajoin_{LR}$}
  \begin{columns}[T]
    \begin{column}{0.48\textwidth}
      \centering
      \vennjoin{0}{1}{1}\\[4pt]
      \small \textbf{Right Outer Join}

      \vspace{4pt}
      \begin{block}{}
        All tuples from the \textbf{right} relation are kept.  Unmatched
        left-side columns become \texttt{NULL}.
      \end{block}

      \small
      \textbf{Products} $\rajoin_R$ \textbf{Brands}:
      \begin{tabular}{llll}
        \toprule
        ProdID & Name & BrandID & BrandName\\
        \midrule
        P1 & Headphones & B1 & Sony \\
        P2 & Speaker    & B2 & Bose \\
        \texttt{NULL} & \texttt{NULL} & B3 & JBL \\
        \bottomrule
      \end{tabular}
    \end{column}
    \begin{column}{0.48\textwidth}
      \centering
      \vennjoin{1}{1}{1}\\[4pt]
      \small \textbf{Full Outer Join}

      \vspace{4pt}
      \begin{block}{}
        \textbf{All} tuples from both sides; \texttt{NULL}-filled on the
        missing side.
        \[
          R_1 \rajoin_{LR} R_2 = (R_1 \rajoin_L R_2) \cup (R_1 \rajoin_R R_2)
        \]
      \end{block}

      \small
      \textbf{Products} $\rajoin_{LR}$ \textbf{Brands}:
      \begin{tabular}{llll}
        \toprule
        ProdID & Name & BrandID & BrandName\\
        \midrule
        P1 & Headphones & B1 & Sony \\
        P2 & Speaker    & B2 & Bose \\
        P3 & Cable      & \texttt{NULL} & \texttt{NULL} \\
        \texttt{NULL} & \texttt{NULL} & B3 & JBL \\
        \bottomrule
      \end{tabular}
    \end{column}
  \end{columns}
\end{frame}

% ----------------------------------------------------------
\begin{frame}{Anti Join $\triangleright$}
  The anti join returns tuples from $R_1$ for which there is \textbf{no}
  matching tuple in $R_2$ --- the complement of a semi-join.

  \begin{columns}[c]
    \begin{column}{0.3\textwidth}
      \centering
      \vennjoin{1}{0}{0}\\[4pt]
      \small Only the \textbf{left-only} part.
    \end{column}
    \begin{column}{0.65\textwidth}
      \[
        X \;\triangleright_{X.A\,=\,Y.A}\; Y
      \]

      \begin{examples}{Use Case --- ``Customers who have never ordered''}
        \[
          \texttt{Customers} \;\triangleright_{\texttt{ID}\,=\,\texttt{CustomerID}}\; \texttt{Orders}
        \]
      \end{examples}

      \begin{examples}{Concrete example (Products $\triangleright$ Brands)}
        Only P3 appears in the result because it is the only product with no
        matching brand.
      \end{examples}

      \vspace{4pt}
      \textbf{Note:} The anti join can be expressed as:\\
      $R_1 \triangleright R_2 = R_1 \setminus (R_1 \ralsem R_2)$
    \end{column}
  \end{columns}
\end{frame}

% ----------------------------------------------------------
\begin{frame}{Join Summary --- All Variants Side by Side}
  \begin{center}
    \small
    \begin{tabular}{llc}
      \toprule
      \textbf{Join Type} & \textbf{Symbol} & \textbf{Venn}\\
      \midrule
      Inner / Theta Join    & $\rajoin_C$      & \vennjoin{0}{1}{0} \\
      Left Outer Join       & $\rajoin_L$      & \vennjoin{1}{1}{0} \\
      Right Outer Join      & $\rajoin_R$      & \vennjoin{0}{1}{1} \\
      Full Outer Join       & $\rajoin_{LR}$   & \vennjoin{1}{1}{1} \\
      Left Semi-Join        & $\ralsem$        & \vennjoin{1}{1}{0} \\
      Anti Join             & $\triangleright$ & \vennjoin{1}{0}{0} \\
      \bottomrule
    \end{tabular}
  \end{center}

  \begin{block}{Key Insight}
    \textbf{Inner:} only overlap. \quad
    \textbf{Full Outer:} everything. \quad
    \textbf{Anti:} orphans only.
  \end{block}
\end{frame}
